\section{TRUY VẤN BẰNG SQL}
\label{sec:chuong5}

Chương này trình bày mười truy vấn QSQL01--QSQL10 trên cơ sở dữ liệu gồm 636 bản ghi. Mỗi câu gồm yêu cầu nghiệp vụ, câu SQL và nhận xét về đầu ra. Mã SQL được lấy từ \texttt{sql/04\_queries.sql} và giữ đúng điều kiện lọc của lần kiểm thử trên pgAdmin. Các ảnh thực thi được dẫn trong từng câu và tập hợp tại \hyperref[sec:anhqsql]{phụ lục ảnh SQL}.\par\medskip

\noindent\begin{minipage}{\linewidth}
\subsection*{QSQL01 --- Cư dân đang cư trú tại Tổ 01}

Liệt kê mã, họ tên và loại cư trú của những người có đợt cư trú còn hiệu lực tại Tổ 01.

\begin{lstlisting}[style=sqlvn]
SELECT DISTINCT
    c.ma_cu_dan,
    c.ho_ten,
    q.loai_cu_tru
FROM cu_dan AS c
JOIN qua_trinh_cu_tru AS q ON q.ma_cu_dan = c.ma_cu_dan
JOIN ho_cu_tru AS h        ON h.ma_ho = q.ma_ho
JOIN dia_chi AS d          ON d.ma_dia_chi = h.ma_dia_chi
JOIN to_dan_pho AS t       ON t.ma_to = d.ma_to
WHERE t.ten_to = 'Tổ 01'
  AND q.ngay_bat_dau <= DATE '2026-11-20'
  AND (q.ngay_ket_thuc IS NULL OR q.ngay_ket_thuc >= DATE '2026-11-20')
ORDER BY c.ma_cu_dan, q.loai_cu_tru;
\end{lstlisting}

Kết quả có 36 dòng thuộc 34 cư dân. \texttt{DISTINCT} loại dòng trùng theo mã cư dân, họ tên và loại cư trú. Một cư dân có cả thường trú và tạm trú còn hiệu lực tại tổ sẽ xuất hiện ở hai dòng với hai loại khác nhau. Xem Hình~\ref{fig:qsql01}.
\end{minipage}\par\medskip

\noindent\begin{minipage}{\linewidth}
\subsection*{QSQL02 --- Hộ có ít nhất 4 cư dân}

Tìm mã các hộ có từ bốn cư dân phân biệt đang cư trú và đếm số cư dân của từng hộ.

\begin{lstlisting}[style=sqlvn]
SELECT
    h.ma_ho,
    COUNT(DISTINCT q.ma_cu_dan) AS so_cu_dan
FROM ho_cu_tru AS h
JOIN qua_trinh_cu_tru AS q ON q.ma_ho = h.ma_ho
WHERE q.ngay_bat_dau <= DATE '2026-11-20'
  AND (q.ngay_ket_thuc IS NULL OR q.ngay_ket_thuc >= DATE '2026-11-20')
GROUP BY h.ma_ho
HAVING COUNT(DISTINCT q.ma_cu_dan) >= 4
ORDER BY so_cu_dan DESC, h.ma_ho;
\end{lstlisting}

Có 24 hộ thỏa điều kiện, nhiều nhất là 7 cư dân. \texttt{WHERE} lọc đợt còn hiệu lực trước khi gom nhóm, còn \texttt{HAVING} giữ các nhóm đạt ngưỡng 4 người. \texttt{COUNT(DISTINCT ...)} đếm một người một lần trong mỗi hộ. Xem Hình~\ref{fig:qsql02}.
\end{minipage}\par\medskip

\noindent\begin{minipage}{\linewidth}
\subsection*{QSQL03 --- Thống kê cư dân theo tổ dân phố}

Đếm số cư dân đang cư trú của từng tổ, kể cả tổ chưa có người.

\begin{lstlisting}[style=sqlvn]
SELECT
    t.ma_to,
    t.ten_to,
    COUNT(DISTINCT q.ma_cu_dan) AS so_cu_dan
FROM to_dan_pho AS t
LEFT JOIN dia_chi AS d          ON d.ma_to = t.ma_to
LEFT JOIN ho_cu_tru AS h        ON h.ma_dia_chi = d.ma_dia_chi
LEFT JOIN qua_trinh_cu_tru AS q ON q.ma_ho = h.ma_ho
                               AND q.ngay_bat_dau <= DATE '2026-11-20'
                               AND (q.ngay_ket_thuc IS NULL OR q.ngay_ket_thuc >= DATE '2026-11-20')
GROUP BY t.ma_to, t.ten_to
ORDER BY t.ma_to;
\end{lstlisting}

Kết quả có 8 dòng, mỗi tổ một dòng. Tổ 05 có 41 cư dân và Tổ 08 có 0 cư dân. Truy vấn bắt đầu từ bảng tổ và đặt điều kiện hiệu lực trong \texttt{ON} của phép nối trái để giữ tổ trống. Số đếm dùng mã cư dân phân biệt trên cả hai loại cư trú. Xem Hình~\ref{fig:qsql03}.
\end{minipage}\par\medskip

\noindent\begin{minipage}{\linewidth}
\subsection*{QSQL04 --- Địa chỉ chưa có hộ cư trú}

Liệt kê mã và địa chỉ chuẩn của những địa chỉ chưa gắn với hộ nào.

\begin{lstlisting}[style=sqlvn]
SELECT
    d.ma_dia_chi,
    d.dia_chi_chuan
FROM dia_chi AS d
WHERE NOT EXISTS (
    SELECT 1
    FROM ho_cu_tru AS h
    WHERE h.ma_dia_chi = d.ma_dia_chi
)
ORDER BY d.ma_dia_chi;
\end{lstlisting}

Có 4 địa chỉ thỏa mãn, với các mã 47 đến 50. \texttt{NOT EXISTS} kiểm tra trực tiếp rằng địa chỉ đang xét không có dòng hộ nào tham chiếu tới nó. Xem Hình~\ref{fig:qsql04}.
\end{minipage}\par\medskip

\noindent\begin{minipage}{\linewidth}
\subsection*{QSQL05 --- Cư dân từng ở ít nhất hai hộ}

Tìm những người có các đợt cư trú đã ghi nhận tại ít nhất hai hộ khác nhau và đếm số hộ theo từng người.

\begin{lstlisting}[style=sqlvn]
SELECT
    c.ma_cu_dan,
    c.ho_ten,
    COUNT(DISTINCT q.ma_ho) AS so_ho
FROM cu_dan AS c
JOIN qua_trinh_cu_tru AS q ON q.ma_cu_dan = c.ma_cu_dan
GROUP BY c.ma_cu_dan, c.ho_ten
HAVING COUNT(DISTINCT q.ma_ho) >= 2
ORDER BY so_ho DESC, c.ma_cu_dan;
\end{lstlisting}

Có 54 cư dân có đợt cư trú được ghi nhận ở từ hai hộ. \texttt{COUNT(DISTINCT q.ma\_ho)} đếm mỗi hộ một lần dù người đó từng có nhiều đợt tại cùng hộ. Chẳng hạn, cư dân mã 92 có ba hộ trong kết quả. Xem Hình~\ref{fig:qsql05}.
\end{minipage}\par\medskip

\noindent\begin{minipage}{\linewidth}
\subsection*{QSQL06 --- Cư dân có cả thường trú và tạm trú}

Tìm những cư dân có ít nhất một đợt thường trú và một đợt tạm trú trong lịch sử đã ghi nhận.

\begin{lstlisting}[style=sqlvn]
SELECT
    c.ma_cu_dan,
    c.ho_ten
FROM cu_dan AS c
WHERE EXISTS (
    SELECT 1
    FROM qua_trinh_cu_tru AS q1
    WHERE q1.ma_cu_dan = c.ma_cu_dan
      AND q1.loai_cu_tru = 'THUONG_TRU'
)
  AND EXISTS (
    SELECT 1
    FROM qua_trinh_cu_tru AS q2
    WHERE q2.ma_cu_dan = c.ma_cu_dan
      AND q2.loai_cu_tru = 'TAM_TRU'
)
ORDER BY c.ma_cu_dan;
\end{lstlisting}

Có 31 cư dân thỏa đồng thời hai điều kiện \texttt{EXISTS}. Hai điều kiện xét độc lập từng loại đợt trong lịch sử, vì vậy hai đợt có thể diễn ra ở những thời điểm khác nhau. Kiểm chứng theo mã cư dân khớp với phép giao ở Câu 5 của Chương \ref{sec:chuong4}. Xem Hình~\ref{fig:qsql06}.
\end{minipage}\par\medskip

\noindent\begin{minipage}{\linewidth}
\subsection*{QSQL07 --- Cư dân từng cư trú và không còn đợt hiệu lực}

Tìm những người có đợt cư trú đã bắt đầu nhưng không còn đợt nào hiệu lực theo điều kiện trong truy vấn.

\begin{lstlisting}[style=sqlvn]
SELECT
    c.ma_cu_dan,
    c.ho_ten
FROM cu_dan AS c
WHERE EXISTS (
    SELECT 1
    FROM qua_trinh_cu_tru AS q_past
    WHERE q_past.ma_cu_dan = c.ma_cu_dan
      AND q_past.ngay_bat_dau <= DATE '2026-11-20'
)
  AND NOT EXISTS (
    SELECT 1
    FROM qua_trinh_cu_tru AS q_cur
    WHERE q_cur.ma_cu_dan = c.ma_cu_dan
      AND q_cur.ngay_bat_dau <= DATE '2026-11-20'
      AND (q_cur.ngay_ket_thuc IS NULL OR q_cur.ngay_ket_thuc >= DATE '2026-11-20')
)
ORDER BY c.ma_cu_dan;
\end{lstlisting}

Có 17 cư dân từng có đợt đã bắt đầu nhưng không còn đợt hiệu lực. \texttt{EXISTS} yêu cầu ít nhất một đợt đã bắt đầu và \texttt{NOT EXISTS} loại cư dân còn đợt hiệu lực. Điều kiện thứ nhất loại cả người chỉ có đợt bắt đầu trong tương lai. Xem Hình~\ref{fig:qsql07}.
\end{minipage}\par\medskip

\noindent\begin{minipage}{\linewidth}
\subsection*{QSQL08 --- Thống kê biến động theo tháng và loại}

Đếm số biến động theo từng cặp tháng và loại biến động trong toàn bộ dữ liệu đã ghi.

\begin{lstlisting}[style=sqlvn]
SELECT
    TO_CHAR(DATE_TRUNC('month', b.ngay_bien_dong), 'YYYY-MM') AS thang,
    b.loai_bien_dong,
    COUNT(*) AS so_su_kien
FROM bien_dong_cu_tru AS b
GROUP BY DATE_TRUNC('month', b.ngay_bien_dong), b.loai_bien_dong
ORDER BY thang, b.loai_bien_dong;
\end{lstlisting}

Kết quả có 45 nhóm tháng và loại, với tổng cộng 63 biến động. Truy vấn gom theo tháng lịch rồi hiển thị theo định dạng \texttt{YYYY-MM}. Bảng dưới trích ba dòng đầu và hiển thị đầy đủ giá trị tháng. Xem Hình~\ref{fig:qsql08}.

\begin{center}
\begin{tabular}{lll}
\hline
\textbf{Tháng} & \textbf{Loại biến động} & \textbf{Số sự kiện} \\
\hline
2025-01 & \texttt{CHUYEN\_DEN} & 1 \\
2025-01 & \texttt{CHUYEN\_DI} & 2 \\
2025-01 & \texttt{CHUYEN\_NOI\_BO} & 1 \\
\hline
\end{tabular}
\end{center}
\end{minipage}\par\medskip

\noindent\begin{minipage}{\linewidth}
\subsection*{QSQL09 --- Chuyển nội bộ giữa hai tổ}

Liệt kê mã biến động, mã cư dân, tổ cũ, tổ mới và ngày của mỗi lần chuyển nội bộ giữa hai tổ khác nhau.

\begin{lstlisting}[style=sqlvn]
SELECT
    b.ma_bien_dong,
    q_cu.ma_cu_dan,
    d_cu.ma_to AS ma_to_cu,
    d_moi.ma_to AS ma_to_moi,
    b.ngay_bien_dong
FROM bien_dong_cu_tru AS b
JOIN qua_trinh_cu_tru AS q_cu  ON q_cu.ma_qua_trinh = b.ma_qua_trinh_cu
JOIN ho_cu_tru AS h_cu         ON h_cu.ma_ho = q_cu.ma_ho
JOIN dia_chi AS d_cu           ON d_cu.ma_dia_chi = h_cu.ma_dia_chi
JOIN qua_trinh_cu_tru AS q_moi ON q_moi.ma_qua_trinh = b.ma_qua_trinh_moi
JOIN ho_cu_tru AS h_moi        ON h_moi.ma_ho = q_moi.ma_ho
JOIN dia_chi AS d_moi          ON d_moi.ma_dia_chi = h_moi.ma_dia_chi
WHERE b.loai_bien_dong = 'CHUYEN_NOI_BO'
  AND q_cu.ma_cu_dan = q_moi.ma_cu_dan
  AND d_cu.ma_to <> d_moi.ma_to
ORDER BY b.ma_bien_dong;
\end{lstlisting}

Kết quả gồm 22 lần chuyển giữa hai tổ. Hai nhánh nối đi từ đợt cũ và đợt mới qua hộ đến địa chỉ để lấy mã tổ. Điều kiện kiểm tra cùng cư dân và tổ cũ khác tổ mới. Chẳng hạn, biến động 2 của cư dân mã 15 là lần chuyển từ Tổ 01 sang Tổ 02 ngày 15/10/2026. Xem Hình~\ref{fig:qsql09}.
\end{minipage}\par\medskip

\noindent\begin{minipage}{\linewidth}
\subsection*{QSQL10 --- Số biến động theo tài khoản}

Đếm số biến động do từng tài khoản ghi nhận, kể cả tài khoản chưa ghi lần nào.

\begin{lstlisting}[style=sqlvn]
SELECT
    a.ma_tai_khoan,
    a.ten_hien_thi,
    COUNT(b.ma_bien_dong) AS so_su_kien
FROM tai_khoan AS a
LEFT JOIN bien_dong_cu_tru AS b ON b.ma_tai_khoan = a.ma_tai_khoan
GROUP BY a.ma_tai_khoan, a.ten_hien_thi
ORDER BY a.ma_tai_khoan;
\end{lstlisting}

Kết quả có 5 tài khoản và tổng 63 biến động. Tài khoản mã 2 ghi 21 biến động, còn tài khoản mã 5 có 0. \texttt{LEFT JOIN} giữ tài khoản chưa ghi và \texttt{COUNT(b.ma\_bien\_dong)} bỏ qua mã biến động \texttt{NULL} ở dòng nối. Xem Hình~\ref{fig:qsql10}.
\end{minipage}\par\medskip

\noindent\begin{minipage}{\linewidth}
\subsection*{Tổng hợp kết quả}

Nhóm chạy lại cả mười truy vấn trên cơ sở dữ liệu dựng mới và kiểm tra đầy đủ các giá trị đầu ra với dữ liệu đối chiếu đã lưu.

\par\medskip
{
\centering
\fontsize{11pt}{14pt}\selectfont
\renewcommand{\arraystretch}{1.15}
\begin{tabularx}{\textwidth}{|P{1.7cm}|Y|P{1.8cm}|}
\hline
\textbf{Mã câu} & \textbf{Kết quả truy vấn} & \textbf{Số dòng} \\
\hline
QSQL01 & Cư dân đang cư trú tại Tổ 01 & 36 \\ \hline
QSQL02 & Hộ có ít nhất 4 cư dân & 24 \\ \hline
QSQL03 & Thống kê cư dân theo tổ dân phố & 8 \\ \hline
QSQL04 & Địa chỉ chưa có hộ cư trú & 4 \\ \hline
QSQL05 & Cư dân từng ở ít nhất hai hộ & 54 \\ \hline
QSQL06 & Cư dân có cả thường trú và tạm trú & 31 \\ \hline
QSQL07 & Cư dân từng cư trú và không còn đợt hiệu lực & 17 \\ \hline
QSQL08 & Thống kê biến động theo tháng và loại & 45 \\ \hline
QSQL09 & Chuyển nội bộ giữa hai tổ & 22 \\ \hline
QSQL10 & Số biến động theo tài khoản & 5 \\ \hline
\end{tabularx}
\captionof{table}{Kết quả kiểm thử QSQL01--QSQL10 trên bộ 636 bản ghi}
\label{tab:ketquasql}
}
\end{minipage}\par\medskip
